\subsection{Competencias específicas, criterios y descriptores}
\vspace{4pt}

\begin{tablaSAlarga}{colspec={Q[l,wd=8.2cm] Q[l,wd=1.5cm] Q[l,wd=3.7cm] Q[l,wd=9.6cm]},
  rowsep=2pt, colsep=4pt, cells={font=\scriptsize}, row{1}={bg=urjc, fg=white, font=\footnotesize\bfseries},}
Competencia específica & Criterios & Descriptores del decreto (en negrita, los priorizados) & Descriptores no priorizados y motivo \\
\textbf{CE1.} Modelizar y resolver problemas de la vida cotidiana y de la ciencia y la tecnología aplicando diferentes estrategias y formas de razonamiento para obtener posibles soluciones. & 1.1, 1.2 & STEM1, STEM2, \textbf{STEM3}, CD2, \textbf{CD5}, CPSAA4, CPSAA5, \textbf{CE3} & STEM1, STEM2 y CD2 se concretan en CE2, CE3 y CE7. CPSAA4 y CPSAA5 no tienen aquí una actividad específica. \\
\textbf{CE2.} Verificar la validez de las posibles soluciones de un problema empleando el razonamiento y la argumentación para contrastar su idoneidad. & 2.1, 2.2 & \textbf{STEM1}, \textbf{STEM2}, CD3, CPSAA4, CC3, CE3 & CPSAA4 se refiere a analizar información de medios de comunicación, y cotejar un cálculo no lo es. CD3, CC3 y CE3 no intervienen en la verificación. \\
\textbf{CE3.} Formular o investigar conjeturas o problemas, utilizando el razonamiento, la argumentación, la creatividad y el uso de herramientas tecnológicas, para generar nuevo conocimiento matemático. & 3.1, 3.2 & CCL1, \textbf{STEM1}, \textbf{STEM2}, CD1, CD2, CD3, CD5, CE3 & CD5 no se moviliza: al conjeturar, GeoGebra es un instrumento de investigación y no se crea ninguna herramienta. CD1, CD2, CD3, CE3 y CCL1 no definen lo que se aprende en esta tarea. \\
\textbf{CE5.} Establecer, investigar y utilizar conexiones entre las diferentes ideas matemáticas estableciendo vínculos entre conceptos, procedimientos, argumentos y modelos para dar significado y estructurar el aprendizaje matemático. & 5.1, 5.2 & \textbf{STEM1}, STEM3, CD2, CD3, CCEC1 & Compartir una construcción (CD3) o crear contenido digital (CD2) no conecta ideas matemáticas. STEM3 se concreta en CE1 y CCEC1 no tiene actividad. \\
\textbf{CE6.} Descubrir los vínculos de las matemáticas con otras áreas de conocimiento y profundizar en sus conexiones, interrelacionando conceptos y procedimientos, para modelizar, resolver problemas y desarrollar la capacidad crítica, creativa e innovadora en situaciones diversas. & 6.1, 6.2 & STEM1, \textbf{STEM2}, CD2, CPSAA5, \textbf{CC4}, CE2, CE3, CCEC1 & CE2 exigiría un análisis económico que la propuesta no incluye. CCEC1 se refiere al patrimonio cultural, y Eratóstenes solo se usa como motivación. STEM1, CD2 y CE3 se concretan en otras competencias; CPSAA5 no tiene actividad. \\
\textbf{CE7.} Representar conceptos, procedimientos e información matemáticos seleccionando diferentes tecnologías, para visualizar ideas y estructurar razonamientos matemáticos. & 7.1, 7.2 & STEM3, CD1, \textbf{CD2}, CD5, CE3, CCEC4.1, CCEC4.2 & CCEC4.1 y CCEC4.2 se refieren a proyectos artísticos, no al diseño de una construcción. CD5, STEM3 y CE3 se concretan una sola vez, en CE1. CD1 no tiene actividad específica. \\
\textbf{CE8.} Comunicar las ideas matemáticas, empleando el soporte, la terminología y el rigor apropiados, para organizar y consolidar el pensamiento matemático. & 8.1 & \textbf{CCL1}, CCL3, CP1, STEM2, \textbf{STEM4}, \textbf{CD3}, CCEC3.2 & CCL3 y CCEC3.2 no tienen actividad específica. CP1 no se evalúa, porque la lengua de trabajo con el socio no es objeto de esta materia. STEM2 se concreta en CE2 y CE3. \\
\textbf{CE9.} Utilizar destrezas personales y sociales, y organizando activamente el trabajo en equipos heterogéneos, aprendiendo del error como parte del proceso de aprendizaje y afrontando situaciones de incertidumbre, para perseverar en la consecución de objetivos en el aprendizaje de las matemáticas. & 9.1, 9.2, 9.3 & CP3, STEM5, \textbf{CPSAA1.1}, CPSAA1.2, CPSAA3.1, \textbf{CPSAA3.2}, CC2, CC3, CE2 & Los demás descriptores no tienen una actividad y una evidencia propias en esta secuencia. \\
\end{tablaSAlarga}

\subsection{Descriptores priorizados: actuación y justificación}
\vspace{4pt}

\begin{tablaSAlarga}{colspec={Q[l,wd=1cm] Q[l,wd=5.6cm] Q[l,wd=5.8cm] Q[l,wd=6.6cm] Q[l,wd=4cm]},
  rowsep=2pt, colsep=4pt, cells={font=\scriptsize}, row{1}={bg=urjc, fg=white, font=\footnotesize\bfseries},}
CE & Descriptor (fragmento literal) & Actuación concreta & Por qué la conexión es real & Evidencia \\
\textbf{CE1} & \textbf{STEM3.} «diseñando y creando prototipos o modelos para generar o utilizar productos que den solución a una necesidad o problema de forma colaborativa» & Cada equipo diseña el modelo de sombras y, con él, la propuesta de ampliación. & El reto es un proyecto: hay una necesidad real del centro y un producto que se evalúa según sus objetivos y su sostenibilidad. & Libro y propuesta (R1), S14--S16 \\
\textbf{CE1} & \textbf{CD5.} «Desarrolla soluciones tecnológicas innovadoras y sostenibles para dar respuesta a necesidades concretas» & El modelo de GeoGebra se entrega como herramienta que la dirección puede reutilizar cambiando la fecha y la altura del obstáculo. & Aquí sí se crea una herramienta digital que responde a una necesidad concreta. Por eso CD5 se trabaja en esta competencia y en ninguna otra. & Construcción dinámica del libro (R1, modelo y método), S15 \\
\textbf{CE1} & \textbf{CE3.} «estrategias ágiles de planificación y gestión de proyectos» & Cada equipo planifica el producto con un tablero Kanban y lo revisa en cada sesión de trabajo. & Kanban es una estrategia ágil, y S14--S17 obligan a gestionar tareas y tiempo para llegar a un resultado útil para otros. & Tablero Kanban (L1), S2--S17 \\
\textbf{CE2} & \textbf{STEM1.} «analizando críticamente las soluciones y reformulando el procedimiento» & El alumnado compara el cálculo a mano con GeoGebra, explica la discrepancia y rehace el procedimiento si hace falta. & El descriptor describe lo mismo que pide la competencia: analizar una solución y corregir el camino seguido. & F1, F3 y examen, pregunta 2 \\
\textbf{CE2} & \textbf{STEM2.} «apreciando la importancia de la precisión y la veracidad» & Cada medida se repite y se calcula cómo cambia la altura con un grado de error. & A diferencia de un problema de libro, la medida real tiene error, y de él depende la validez de la solución. & Análisis del error en el libro (R1, error y solución), S14--S15 \\
\textbf{CE3} & \textbf{STEM1.} «métodos inductivos y deductivos propios del razonamiento matemático» & En F2 se observan casos y se formula una conjetura (S5); después se demuestra (S6). & La secuencia reproduce los dos métodos que nombra el descriptor: primero inducción y luego deducción. & F2 y examen, pregunta 3 \\
\textbf{CE3} & \textbf{STEM2.} «planteándose hipótesis y contrastándolas o comprobándolas mediante la observación» & En F2 se pone a prueba en GeoGebra la conjetura $\sin 2\alpha = 2\sin\alpha$ y se busca un contraejemplo. & La conjetura funciona como hipótesis y GeoGebra como instrumento de observación, no como producto. & F2 y observación (L1), S5 \\
\textbf{CE5} & \textbf{STEM1.} «emplea estrategias variadas para la resolución de problemas» & Ángulos, perpendicularidad y distancias se abordan con producto escalar, ecuaciones de la recta y trigonometría. & Conectar ideas se demuestra al poder resolver un problema por más de una vía y elegir entre ellas. & Mapa conceptual (L2) y examen, pregunta 5 \\
\textbf{CE6} & \textbf{CC4.} «actividades y hábitos que conduzcan al logro de los Objetivos de Desarrollo Sostenible» & Cada equipo argumenta cómo contribuye su propuesta a los ODS 7 y 11. & La decisión afecta a la energía de su propio centro: una actuación local conectada con un objetivo global. & Reflexión sobre los ODS en el libro (R1, ODS y límites) y defensa (R2) \\
\textbf{CE6} & \textbf{STEM2.} «mostrando una actitud crítica acerca del alcance y limitaciones de los métodos empleados» & Se compara la separación del mediodía con la del pliego del IDAE y se señalan los límites del modelo. & El modelo matemático por sí solo no basta: el criterio técnico muestra hasta dónde llega y dónde se queda corto. & Propuesta y límites del modelo en el libro (R1, propuesta y ODS y límites) \\
\textbf{CE7} & \textbf{CD2.} «Crea, integra y reelabora contenidos digitales de forma individual o colectiva» & El equipo integra en el libro construcciones, cálculos y explicación, y cita los datos y materiales ajenos. & El libro es un contenido digital colectivo que reelabora datos de terceros, por lo que también obliga a respetar su autoría. & Construcciones y fuentes citadas en el libro (R1, representación) \\
\textbf{CE8} & \textbf{CCL1.} «adecuación a los diferentes contextos sociales y académicos» & La propuesta se defiende ante la clase y la dirección. & Explicar el mismo método a compañeros y a quien decide obliga a adaptar el discurso sin perder rigor. & Exposición (R2), S19--S20 \\
\textbf{CE8} & \textbf{STEM4.} «de forma clara y precisa, en diferentes formatos» & El método y los resultados se presentan con fórmulas, tablas, gráficos y construcciones. & Cada formato aporta algo distinto: el plano convence a la dirección y la fórmula justifica el cálculo. & Exposición (R2) y libro (R1) \\
\textbf{CE8} & \textbf{CD3.} «para comunicarse, trabajar colaborativamente y compartir información, gestionando de manera responsable sus acciones, presencia y visibilidad en la red» & Datos y conclusiones se comparten en el TwinSpace con avatares y sin datos personales. & Es un intercambio real con otro centro, en el que la identidad digital se gestiona con reglas explícitas. & Publicaciones del TwinSpace (seguimiento, sin calificación) \\
\textbf{CE9} & \textbf{CPSAA3.2.} «Distribuye en un grupo las tareas, recursos y responsabilidades de manera ecuánime» & Los cuatro roles rotan, de modo que todos los ejercen todos\soloProducto{ (sección~\ref{sec:cooperativo})}. & La rotación convierte el reparto ecuánime en una regla de funcionamiento, no en una intención. & Tablero Kanban y cumplimiento del rol (L1), R3 \\
\textbf{CE9} & \textbf{CPSAA1.1.} «la resiliencia, la autoeficacia y la búsqueda de objetivos de forma autónoma» & Las conjeturas falsas y las medidas fallidas se revisan y se repiten. & El reto incluye fallos reales (nubes, medidas incoherentes), que dan ocasión de perseverar. & Observación (L1) y autoevaluación (R3) \\
\end{tablaSAlarga}

\subsection{Criterios operativizados, saberes, sesiones, instrumentos y evidencias}
\vspace{4pt}

\begin{tablaSAlarga}{colspec={Q[l,wd=5.2cm] Q[l,wd=7.2cm] Q[c,wd=1.3cm] Q[c,wd=1.3cm] Q[c,wd=1.5cm] Q[l,wd=6.5cm]},
  rowsep=2pt, colsep=4pt, cells={font=\scriptsize}, row{1}={bg=urjc, fg=white, font=\footnotesize\bfseries},}
Criterio de evaluación & Criterio operativizado & Bloom & Sesiones & Instrum. & Evidencia de aprendizaje \\
\SetCell[c=6]{l,bg=rosa} {\textbf{A. Números y operaciones: vectores}\\ \textbullet~Adición y producto escalar de vectores: propiedades y representaciones.\\ \textbullet~Conjunto de vectores: estructura, comprensión y propiedades.\\ \textbullet~Módulo de un vector, coordenadas de un vector con respecto a una base, ángulo entre dos vectores y proyección ortogonal.\\ \textbullet~Estrategias para operar con números reales y vectores: cálculo mental o escrito en los casos sencillos y con herramientas tecnológicas en los casos más complicados.} & & & & & \\*
\textbf{7.1.} Representar ideas matemáticas, estructurando diferentes razonamientos matemáticos y seleccionando las tecnologías más adecuadas. & El alumno ha representado vectores en GeoGebra y ha calculado su suma, su producto por un número y su producto escalar, interpretando geométricamente cada resultado. & \textit{Aplicar} & 2, 3, 4 & CU, FI & Ficha F1 y cuestionario Q1 con operaciones interpretadas. \\
\textbf{2.1.} Comprobar la validez matemática de las posibles soluciones de un problema utilizando el razonamiento y la argumentación. & El alumno ha comprobado con GeoGebra el módulo, el ángulo entre dos vectores y la proyección ortogonal calculados a mano, justificando el origen de las discrepancias. & \textit{Evaluar} & 2, 4, 18 & EX, FI & Ficha F1 con cálculo manual, comprobación y causa de la discrepancia. \\
\textbf{5.2.} Resolver problemas en contextos matemáticos estableciendo y aplicando conexiones entre las diferentes ideas matemáticas. & El alumno ha resuelto problemas de perpendicularidad y de cálculo de ángulos aplicando el producto escalar y relacionando su expresión en coordenadas con su significado geométrico. & \textit{Aplicar} & 4, 18 & EX, CU & Problemas de perpendicularidad y ángulos resueltos. \\
\SetCell[c=6]{l,bg=rosa} {\textbf{B y D. Trigonometría}\\ \textbullet~Razones trigonométricas de un ángulo cualquiera medido en grados o en radianes.\\ \textbullet~Demostración de las identidades trigonométricas. Razones trigonométricas del ángulo suma, el ángulo diferencia, el ángulo doble y el ángulo mitad.\\ \textbullet~Cálculo de las razones trigonométricas de un ángulo cualquiera empleando las principales fórmulas trigonométricas.\\ \textbullet~Conjeturas geométricas en el plano: validación por medio de la deducción y la demostración de teoremas.\\ \textbullet~Resolución de ecuaciones trigonométricas sencillas.} & & & & & \\*
\textbf{3.2.} Emplear herramientas tecnológicas adecuadas en la formulación o investigación de conjeturas o problemas. & El alumno ha examinado en una circunferencia goniométrica dinámica de GeoGebra (con un deslizador para el ángulo) cómo varían el seno, el coseno y la tangente en los cuatro cuadrantes. & \textit{Analizar} & 5 & FI, OB & Ficha F2 con la tabla de signos y observaciones por cuadrante. \\
\textbf{3.1.} Adquirir nuevo conocimiento matemático a partir de la formulación de conjeturas y problemas de forma guiada. & El alumno ha formulado, de forma guiada, conjeturas sobre identidades trigonométricas (fundamental, ángulo doble, ángulo suma) a partir de la construcción y las ha demostrado por escrito. & \textit{Crear} & 5, 6, 18 & EX, FI & Conjetura escrita en F2 y demostración en papel. \\
\textbf{1.2.} Obtener todas las posibles soluciones matemáticas de problemas de la vida cotidiana y de la ciencia y la tecnología, describiendo el procedimiento utilizado. & El alumno ha resuelto ecuaciones trigonométricas sencillas obteniendo todas sus soluciones, tanto en [0°, 360°) como en forma general, y ha descrito el procedimiento seguido. & \textit{Aplicar} & 7, 18 & EX, CU & Cuestionario Q2: soluciones en el intervalo y solución general. \\
\SetCell[c=6]{l,bg=rosa} {\textbf{B. Medida y geometría: resolución de triángulos}\\ \textbullet~Cálculo de longitudes y medidas angulares: uso de la trigonometría.\\ \textbullet~Aplicación de las razones trigonométricas, el teorema de los senos y el teorema del coseno en la resolución de triángulos y de problemas geométricos de contexto real.\\ \textbullet~Demostración del teorema del seno y del coseno.} & & & & & \\*
\textbf{1.1.} Manejar algunas estrategias y herramientas, incluidas las digitales, en la modelización y resolución de problemas de la vida cotidiana y de la ciencia y la tecnología, evaluando su eficiencia en cada caso. & El alumno ha justificado la elección entre el teorema del seno y el del coseno según los datos de cada triángulo, valorando la eficiencia de cada estrategia. & \textit{Evaluar} & 8, 9, 18 & EX, FI & Ficha F3: elección razonada del teorema. \\
\textbf{1.2.} Obtener todas las posibles soluciones matemáticas de problemas de la vida cotidiana y de la ciencia y la tecnología, describiendo el procedimiento utilizado. & El alumno ha discriminado si un triángulo dado por dos lados y el ángulo opuesto a uno de ellos tiene dos, una o ninguna solución y ha calculado todas las soluciones, describiendo el procedimiento. & \textit{Analizar} & 9, 18 & EX, FI & Ficha F3: número de soluciones y procedimiento. \\
\textbf{2.1.} Comprobar la validez matemática de las posibles soluciones de un problema utilizando el razonamiento y la argumentación. & El alumno ha verificado en GeoGebra, arrastrando los vértices, la validez de los triángulos obtenidos analíticamente y ha argumentado por qué un caso no tiene solución. & \textit{Evaluar} & 9 & CU, FI & Cuestionario Q3 y comprobación arrastrando vértices. \\
\SetCell[c=6]{l,bg=rosa} {\textbf{C. Geometría en el plano: rectas}\\ \textbullet~Ecuaciones de la recta en el espacio bidimensional.\\ \textbullet~Estudio de la posición relativa de puntos y rectas en el plano.\\ \textbullet~Planteamiento y resolución de problemas de geometría afín relacionados con la incidencia, el paralelismo y la ortogonalidad de rectas en el plano.\\ \textbullet~Planteamiento y resolución de problemas de geometría métrica relacionados con la medida de ángulos entre rectas y la medida de distancias entre puntos y rectas.\\ \textbullet~Expresiones algebraicas de objetos geométricos: selección de la más adecuada en función de la situación a resolver.} & & & & & \\*
\textbf{5.1.} Manifestar una visión matemática integrada, investigando y conectando las diferentes ideas matemáticas. & El alumno ha relacionado el vector director y el vector normal con las distintas ecuaciones de la recta, organizando en un mapa conceptual las conexiones entre vectores, trigonometría y geometría analítica. & \textit{Analizar} & 10, 13, 18 & EX, MC & Mapa conceptual con enlaces entre vector director, normal y ángulo. \\
\textbf{5.2.} Resolver problemas en contextos matemáticos estableciendo y aplicando conexiones entre las diferentes ideas matemáticas. & El alumno ha determinado la posición relativa de dos rectas y ha calculado ángulos entre rectas y distancias de un punto a una recta aplicando el producto escalar. & \textit{Aplicar} & 11, 18 & EX, CU & Cuestionario Q4 con posición, ángulo y distancia. \\
\textbf{7.2.} Seleccionar y utilizar diversas formas de representación, valorando su utilidad para compartir información. & El alumno ha seleccionado la ecuación de la recta (vectorial, paramétrica, continua, general o explícita) más útil para cada problema, justificando su elección. & \textit{Evaluar} & 10, 12, 18 & EX, FI & Ficha F4: ecuación elegida y motivo. \\
\SetCell[c=6]{l,bg=rosa} {\textbf{C y D. Modelización: producto final «Medir lo que no se puede alcanzar»}\\ \textbullet~Representación de objetos geométricos en el plano mediante herramientas digitales.\\ \textbullet~Modelos matemáticos (geométricos, algebraicos, grafos...) en la resolución de problemas en el plano. Conexiones con otras disciplinas y áreas de interés.\\ \textbullet~Modelización de la posición y el movimiento de un objeto en el plano mediante vectores.\\ \textbullet~Formulación, resolución y análisis de problemas de la vida cotidiana y de la ciencia y la tecnología utilizando herramientas o programas adecuados.} & & & & & \\*
\textbf{1.1.} Manejar algunas estrategias y herramientas, incluidas las digitales, en la modelización y resolución de problemas de la vida cotidiana y de la ciencia y la tecnología, evaluando su eficiencia en cada caso. & El alumno ha diseñado un método de medida y ha construido un modelo a escala en GeoGebra para calcular una altura o distancia inaccesible del centro, valorando su eficiencia frente a otros métodos. & \textit{Crear} & 14, 15, 19, 20 & LI, EP & Libro: método, construcción y comparación con otro método. \\
\textbf{2.2.} Seleccionar la solución más adecuada de un problema en función del contexto, usando el razonamiento y la argumentación. & El alumno ha valorado el efecto de un error de un grado en el resultado final y ha seleccionado, entre las medidas repetidas, el resultado más razonable. & \textit{Evaluar} & 14, 15, 19, 20 & LI, EP & Libro: medidas repetidas, sensibilidad y resultado elegido. \\
\textbf{7.2.} Seleccionar y utilizar diversas formas de representación, valorando su utilidad para compartir información. & El alumno ha elaborado un libro de GeoGebra que integra la construcción, los cálculos y la explicación del método, eligiendo las representaciones más claras para compartirlo con la clase. & \textit{Crear} & 15, 16, 17, 19, 20 & LI, CO & Libro completo con representaciones justificadas. \\
\textbf{8.1.} Mostrar organización al comunicar las ideas matemáticas empleando el soporte, la terminología y el rigor apropiados. & El alumno ha explicado oralmente el método y los resultados con orden, utilizando la terminología y la notación matemáticas con rigor. & \textit{Comprender} & 17, 19, 20 & EP, CO & Ensayo y defensa oral con terminología y notación correctas. \\
\SetCell[c=6]{l,bg=rosa} {\textbf{F. Actitudes y aprendizaje}\\ \textbullet~Tratamiento del error como elemento movilizador de conocimientos previos adquiridos y generador de oportunidades de aprendizaje en el aula de matemáticas.\\ \textbullet~Reconocimiento y aceptación de diversos planteamientos en la resolución de problemas, transformando los enfoques de los demás en nuevas y mejoradas estrategias propias.\\ \textbullet~Técnicas y estrategias de trabajo en equipo para la resolución de problemas y tareas matemáticas, en grupos heterogéneos.\\ \textbullet~Destrezas para desarrollar una comunicación efectiva: la escucha activa, la formulación de preguntas o solicitud y prestación de ayuda cuando sea necesario.} & & & & & \\*
\textbf{9.1.} Afrontar las situaciones de incertidumbre, aceptando y aprendiendo del error como parte del proceso de aprendizaje de las matemáticas. & El alumno ha analizado los errores de sus conjeturas y de sus medidas, y los ha utilizado como punto de partida para revisar su razonamiento en la puesta en común. & \textit{Analizar} & 5, 12, 14, 19, 20 & OB, CO & Registro de observación y revisión de medidas fallidas. \\
\textbf{9.2.} Mostrar una actitud positiva y perseverante, aceptando y aprendiendo de la crítica razonada al hacer frente a las diferentes actividades de las matemáticas. & El alumno ha valorado la crítica razonada recibida en la coevaluación y la ha incorporado para mejorar el libro de GeoGebra, perseverando ante las dificultades. & \textit{Evaluar} & 16, 17, 19, 20 & OB, CO & Registro de cambios del libro tras la coevaluación. \\
\textbf{9.3.} Participar en tareas matemáticas de forma activa en equipos heterogéneos, escuchando su razonamiento. & El alumno ha colaborado activamente en el equipo heterogéneo, ha asumido tareas del tablero Kanban y ha tenido en cuenta los razonamientos de sus compañeros. & \textit{Aplicar} & 2, 14, 15, 16, 19, 20 & OB, CO & Tablero Kanban, roles rotados y registro de observación. \\
\SetCell[c=6]{l,bg=rosa} {\textbf{C y F. Conexión con la energía solar}\\ \textbullet~Modelos matemáticos (geométricos, algebraicos, grafos...) en la resolución de problemas en el plano. Conexiones con otras disciplinas y áreas de interés.\\ \textbullet~Valoración de la contribución de las matemáticas y el papel de matemáticos a lo largo de la historia en el avance de la ciencia y la tecnología.} & & & & & \\*
\textbf{6.1.} Resolver problemas en situaciones diversas, utilizando procesos matemáticos, estableciendo y aplicando conexiones entre el mundo real, otras áreas de conocimiento y las matemáticas. & El alumno ha resuelto la ubicación de los nuevos paneles conectando las medidas reales, el modelo geométrico y el criterio técnico de separación entre filas. & \textit{Aplicar} & 14, 16, 19, 20 & LI, EP & Propuesta de ampliación con plano, cálculo de filas y decisión justificada. \\
\textbf{6.2.} Analizar la aportación de las matemáticas al progreso de la humanidad reflexionando sobre su contribución en la propuesta de soluciones a situaciones complejas y a los retos científicos y tecnológicos que se plantean en la sociedad. & El alumno ha analizado cómo contribuye su propuesta a los retos de energía sostenible (ODS 7 y 11), señalando el papel del modelo matemático y sus límites. & \textit{Analizar} & 16, 19, 20 & LI, EP & Apartado del libro sobre ODS, límites del modelo y defensa. \\
\end{tablaSAlarga}

\ifjustificacion
\subsection{Justificación pedagógica de las relaciones}
\vspace{4pt}

\begin{tablaSAlarga}{colspec={Q[l,wd=1cm] Q[l,wd=3.4cm] Q[c,wd=1.4cm] Q[l,wd=8.9cm] Q[c,wd=1.3cm] Q[l,wd=7cm]},
  rowsep=2pt, colsep=4pt, cells={font=\scriptsize}, row{1}={bg=urjc, fg=white, font=\footnotesize\bfseries},}
Crit. & Aspecto & Sesiones & Por qué en esas sesiones & Instrum. & Por qué ese instrumento \\
\textbf{7.1} & Representar y operar con vectores & 2, 3, 4 & Se construye primero en F1 (S2), para que cada operación tenga un significado geométrico antes de formalizarla en S3 y S4. & CU, FI & La ficha recoge la representación con tecnología y Q1 comprueba su interpretación. El examen escrito no permite usar GeoGebra. \\
\textbf{2.1} & Validar cálculos vectoriales & 2, 4, 18 & Validar exige tener un cálculo propio: en S2 se contrasta de forma exploratoria y en S4, con la fórmula de la proyección, se explica la discrepancia. S18 comprueba si se hace sin herramienta. & EX, FI & F1 registra el contraste entre mano y GeoGebra; la pregunta 2 del examen obliga a detectar un error ajeno. \\
\textbf{5.2} & Producto escalar en problemas & 4, 18 & El producto escalar une coordenadas y geometría, y se enseña en S4; S18 comprueba que se aplica de forma autónoma. & EX, CU & Q1 da un diagnóstico temprano para el refuerzo por criterio de S13, y el examen verifica la transferencia. \\
\textbf{3.2} & Investigar con la circunferencia & 5 & S5 precede a propósito a la exposición de S6: hay que investigar antes de conocer la fórmula. & FI, OB & La ficha muestra lo investigado; la observación capta si se explora antes de pedir ayuda, que no queda en el papel. \\
\textbf{3.1} & Conjeturar y demostrar & 5, 6, 18 & La conjetura nace en S5 y se demuestra en S6, de modo que la demostración responde a una pregunta propia. S18 comprueba que se demuestra sin apoyo. & EX, FI & F2 recoge la conjetura y la pregunta 3a del examen, la demostración individual. \\
\textbf{1.2} & Ecuaciones trigonométricas & 7, 18 & Las ecuaciones usan el ángulo doble, que se estudia en la misma sesión, y la pieza del reto les da un uso inmediato. & EX, CU & Obtener todas las soluciones es un procedimiento: Q2 da retorno inmediato y el examen lo certifica. \\
\textbf{1.1} & Elegir el teorema adecuado & 8, 9, 18 & La elección se aprende en S8 con datos variados y se pone a prueba en S9, con el caso que genera más dudas. & EX, FI & F3 y la pregunta 4 piden justificar la elección, no solo dar el resultado. \\
\textbf{1.2} & Caso ambiguo & 9, 18 & Va al final del bloque porque exige dominar antes el teorema del seno. & EX, FI & En F3 se ve cuántas soluciones hay al arrastrar; la pregunta 4b exige obtenerlas por cálculo. \\
\textbf{2.1} & Validar triángulos & 9 & S9 concentra el caso en que validar es decisivo: si $\sin B>1$, el triángulo no existe. & CU, FI & La comprobación arrastrando vértices solo se ve en F3; Q3 la cierra con un caso sin solución. \\
\textbf{5.1} & Conectar vector y recta & 10, 13, 18 & S10 deduce las ecuaciones del vector director; el mapa llega en S13, cuando ya se han visto los tres bloques y hay algo que conectar. & EX, MC & El mapa hace visibles las conexiones y la pregunta 5 del examen comprueba que se usan. \\
\textbf{5.2} & Ángulos y distancias & 11, 18 & Reutilizan el producto escalar de S4; por eso se sitúan al final del bloque, para volver sobre él. & EX, CU & Q4 alimenta el refuerzo por criterio de S13 y el examen verifica el dominio individual. \\
\textbf{7.2} & Elegir la ecuación de la recta & 10, 12, 18 & S10 presenta las cinco formas y F4, en S12, obliga a elegir una para un problema real. & EX, FI & F4 y la pregunta 5d piden explicar por qué se elige esa ecuación. \\
\textbf{1.1} & Método y modelo & 14, 15, 19, 20 & Se aplica cuando ya se dispone de todas las herramientas: el método se aprueba en S13, se ejecuta en S14 y se modela en S15. & LI, EP & El libro documenta el proceso; la defensa comprueba que todo el equipo lo entiende. \\
\textbf{2.2} & Elegir el resultado & 14, 15, 19, 20 & Elegir solo tiene sentido con medidas reales y repetidas (S14), y la sensibilidad al error se analiza con ellas en S15. & LI, EP & El libro recoge datos y decisión; en la defensa hay que justificarla ante preguntas. \\
\textbf{7.2} & Integrar representaciones & 15, 16, 17, 19, 20 & Integrar exige tener ya datos y modelo (S15 y S16), y S17 permite mejorarlo; las defensas muestran si otras personas lo entienden. & LI, CO & El libro es la representación. La coevaluación aporta lo que pide el criterio: su utilidad para compartir información. \\
\textbf{8.1} & Comunicar con rigor & 17, 19, 20 & Se ensaya en S17 y se evalúa en S19 y S20, ante un público real. & EP, CO & La comunicación oral solo se puede evaluar en una exposición; la co/autoevaluación añade la mirada del público. \\
\textbf{9.1} & Aprender del error & 5, 12, 14, 19, 20 & Se observa donde el error aparece de forma natural: conjeturas falsas (S5), construcciones (S12) y medidas (S14). & OB, CO & Una actitud no cabe en un examen: se observa cuando ocurre y se reflexiona en la autoevaluación. \\
\textbf{9.2} & Aceptar la crítica & 16, 17, 19, 20 & Necesita una crítica aprovechable: la coevaluación del borrador en S16 se aprovecha en S17, dentro del horario de clase. & OB, CO & El registro de cambios permite ver qué se ha incorporado realmente. \\
\textbf{9.3} & Colaborar en equipo & 2, 14, 15, 16, 19, 20 & Los roles se presentan en S2 y se exigen sobre todo de S14 a S16, cuando el producto depende de todos. & OB, CO & Solo la observación y los compañeros pueden valorar la colaboración efectiva. \\
\textbf{6.1} & Conectar con la instalación & 14, 16, 19, 20 & La conexión con el criterio técnico requiere haber medido; se hace en S16, con la orientación del departamento de Tecnología. & LI, EP & La propuesta recoge el cálculo y la defensa, su viabilidad ante quien decide. \\
\textbf{6.2} & Valorar impacto y límites & 16, 19, 20 & La reflexión tiene sentido cuando el modelo ya existe y se conocen sus resultados. & LI, EP & El apartado del libro la deja por escrito y la defensa la somete a preguntas. \\
\end{tablaSAlarga}

\fi

\clearpage
\subsection{Cobertura: cada criterio por sesiones e instrumentos}
\vspace{4pt}

\begin{tablaSAlarga}{colspec={Q[c,wd=1.3cm] *{20}{Q[c,wd=0.58cm]} *{8}{Q[c,wd=0.75cm]}},
  rowsep=1.5pt, colsep=2pt, cells={font=\scriptsize}, row{1}={bg=urjc, fg=white, font=\scriptsize\bfseries},
  column{22-29}={bg=rosa}, cell{1}{22-29}={bg=urjc},}
Criterio & S1 & S2 & S3 & S4 & S5 & S6 & S7 & S8 & S9 & S10 & S11 & S12 & S13 & S14 & S15 & S16 & S17 & S18 & S19 & S20 & EX & CU & FI & LI & EP & MC & OB & CO \\
\textbf{1.1} &  &  &  &  &  &  &  & \textbf{X} & \textbf{X} &  &  &  &  & \textbf{X} & \textbf{X} &  &  & \textbf{X} & \textbf{X} & \textbf{X} & \textbf{X} &  & \textbf{X} & \textbf{X} & \textbf{X} &  &  &  \\
\textbf{1.2} &  &  &  &  &  &  & \textbf{X} &  & \textbf{X} &  &  &  &  &  &  &  &  & \textbf{X} &  &  & \textbf{X} & \textbf{X} & \textbf{X} &  &  &  &  &  \\
\textbf{2.1} &  & \textbf{X} &  & \textbf{X} &  &  &  &  & \textbf{X} &  &  &  &  &  &  &  &  & \textbf{X} &  &  & \textbf{X} & \textbf{X} & \textbf{X} &  &  &  &  &  \\
\textbf{2.2} &  &  &  &  &  &  &  &  &  &  &  &  &  & \textbf{X} & \textbf{X} &  &  &  & \textbf{X} & \textbf{X} &  &  &  & \textbf{X} & \textbf{X} &  &  &  \\
\textbf{3.1} &  &  &  &  & \textbf{X} & \textbf{X} &  &  &  &  &  &  &  &  &  &  &  & \textbf{X} &  &  & \textbf{X} &  & \textbf{X} &  &  &  &  &  \\
\textbf{3.2} &  &  &  &  & \textbf{X} &  &  &  &  &  &  &  &  &  &  &  &  &  &  &  &  &  & \textbf{X} &  &  &  & \textbf{X} &  \\
\textbf{5.1} &  &  &  &  &  &  &  &  &  & \textbf{X} &  &  & \textbf{X} &  &  &  &  & \textbf{X} &  &  & \textbf{X} &  &  &  &  & \textbf{X} &  &  \\
\textbf{5.2} &  &  &  & \textbf{X} &  &  &  &  &  &  & \textbf{X} &  &  &  &  &  &  & \textbf{X} &  &  & \textbf{X} & \textbf{X} &  &  &  &  &  &  \\
\textbf{6.1} &  &  &  &  &  &  &  &  &  &  &  &  &  & \textbf{X} &  & \textbf{X} &  &  & \textbf{X} & \textbf{X} &  &  &  & \textbf{X} & \textbf{X} &  &  &  \\
\textbf{6.2} &  &  &  &  &  &  &  &  &  &  &  &  &  &  &  & \textbf{X} &  &  & \textbf{X} & \textbf{X} &  &  &  & \textbf{X} & \textbf{X} &  &  &  \\
\textbf{7.1} &  & \textbf{X} & \textbf{X} & \textbf{X} &  &  &  &  &  &  &  &  &  &  &  &  &  &  &  &  &  & \textbf{X} & \textbf{X} &  &  &  &  &  \\
\textbf{7.2} &  &  &  &  &  &  &  &  &  & \textbf{X} &  & \textbf{X} &  &  & \textbf{X} & \textbf{X} & \textbf{X} & \textbf{X} & \textbf{X} & \textbf{X} & \textbf{X} &  & \textbf{X} & \textbf{X} &  &  &  & \textbf{X} \\
\textbf{8.1} &  &  &  &  &  &  &  &  &  &  &  &  &  &  &  &  & \textbf{X} &  & \textbf{X} & \textbf{X} &  &  &  &  & \textbf{X} &  &  & \textbf{X} \\
\textbf{9.1} &  &  &  &  & \textbf{X} &  &  &  &  &  &  & \textbf{X} &  & \textbf{X} &  &  &  &  & \textbf{X} & \textbf{X} &  &  &  &  &  &  & \textbf{X} & \textbf{X} \\
\textbf{9.2} &  &  &  &  &  &  &  &  &  &  &  &  &  &  &  & \textbf{X} & \textbf{X} &  & \textbf{X} & \textbf{X} &  &  &  &  &  &  & \textbf{X} & \textbf{X} \\
\textbf{9.3} &  & \textbf{X} &  &  &  &  &  &  &  &  &  &  &  & \textbf{X} & \textbf{X} & \textbf{X} &  &  & \textbf{X} & \textbf{X} &  &  &  &  &  &  & \textbf{X} & \textbf{X} \\
\end{tablaSAlarga}
